\chapter{Complexity — Dbl Lang}
%%% generated by the blueprint pipeline from TCSlib.Complexity.SpaceComplexity.Machines.DblLang
\section{Overview}

Declarations of the Lean module \texttt{TCSlib.Complexity.SpaceComplexity.Machines.DblLang}.

\section{Declarations}

\begin{definition}[Doubled unary and binary language]
\lean{Complexity.LogProg.dblLang}
\label{Complexity.LogProg.dblLang}
\leanok
\uses{Turing.pairEncode}
Define a language of Boolean words consisting of exactly those words for which there is
a natural number $n$ such that the word is the pair encoding of the unary word of $n$
true symbols and the binary expansion of $2n$.
\end{definition}

\begin{definition}[States for the doubled-language machine]
\lean{Complexity.LogProg.Dbl.St}
\label{Complexity.LogProg.Dbl.St}
\leanok
Introduce a state type with the distinct states $vU0$, $vU1$, $vS$, $vW0$, $vWF$, $vWT$,
$rw1$, $rw2$, $cnt$, $iC$, $iB$, $sw1$, $sw2$, $jK$, $jK2$, $jC$, $yes$, and $no$,
together with states $jRB(b)$, $jI1(b)$, and $jI2(b)$ for each Boolean value $b$.
\end{definition}

\begin{definition}[Transition function for the doubling machine]
\lean{Complexity.LogProg.Dbl.tr}
\label{Complexity.LogProg.Dbl.tr}
\leanok
\uses{Complexity.LogProg.Dbl.St, Complexity.LogProg.backAct, Complexity.LogProg.cmpAct, Complexity.LogProg.incBAct, Complexity.LogProg.incCAct, Complexity.LogProg.retAct, Complexity.LogProg.rewAct, Complexity.LogProg.rw2Act, Complexity.LogProg.skipAct, Complexity.LogProg.valSAct, Complexity.LogProg.valUAct, Complexity.LogProg.valWAct, Complexity.LogProg.xAct, Turing.Action}
The transition function assigns an action to each state of the doubling machine, the
symbol under its input head (or blank), and the symbol under its single work-tape head
(or blank), according to the specified state-dependent rules. The action records the
head movements, any work-tape write or output, and the next state or halting decision.
\end{definition}

\begin{definition}[Doubling program]
\lean{Complexity.LogProg.Dbl.prog}
\label{Complexity.LogProg.Dbl.prog}
\leanok
\uses{Complexity.LogProg.Dbl.St, Complexity.LogProg.Dbl.tr, Complexity.LogProg.RProg}
The program has one work tape, no oracle deciders, and the finite state set of the
doubling machine. Its ordinary transition machine starts in state $vU0$ and uses the
transition function specified for that machine; it has no call nodes in any state.
\end{definition}

\begin{definition}[Oracle function with empty domain]
\lean{Complexity.LogProg.Dbl.o}
\label{Complexity.LogProg.Dbl.o}
\leanok
The function on the empty set of oracle indices and Boolean strings returns a Boolean
value; since there are no oracle indices, it has no inputs in its domain.
\end{definition}

\begin{definition}[Initial work-tape configuration]
\lean{Complexity.LogProg.Dbl.K}
\label{Complexity.LogProg.Dbl.K}
\leanok
\uses{Complexity.LogProg.Dbl.St, Turing.Cfg, Turing.FinTM.bufferTape}
For a Boolean input string $y$, a state $s$, an input-head position $ip$ among the
$|y|+2$ positions, a natural number $v$, and an integer work-head position $q$, let
$K(s,ip,v,q)$ be the configuration with state $s$, input-head position $ip$, empty
output, and its single work tape containing the bit representation of $v$ as a buffer,
with the work head at $q$.
\end{definition}

\begin{lemma}[Configuration update for the doubled machine]
\lean{Complexity.LogProg.Dbl.xCfg_K}
\label{Complexity.LogProg.Dbl.xCfg_K}
\leanok
\uses{Complexity.LogProg.Dbl.K, Complexity.LogProg.Dbl.St, Complexity.LogProg.xCfg}
\difficulty{3}
For any states $s,s'$, valid input-head positions $ip,ip'$ on a fixed Boolean input $y$,
natural number $v$, and integers $q,q'$, begin with a configuration in state $s'$,
input-head position $ip'$, tape storing the binary expansion of $v$, and work-tape head
at $q'$. Replacing its state by $s$, its input-head position by $ip$, and its work-tape
head position by $q$ yields the configuration with those updated state and positions,
the same tape, and empty output.
\end{lemma}

\begin{lemma}[Register Update of a Represented Configuration]
\lean{Complexity.LogProg.Dbl.regCfg_K}
\label{Complexity.LogProg.Dbl.regCfg_K}
\leanok
\uses{Complexity.LogProg.Dbl.K, Complexity.LogProg.Dbl.St, Complexity.LogProg.regCfg, Turing.FinTM.bufferTape}
\difficulty{4}
For any states $s,s'$, input-head position $ip$, natural numbers $v,v'$, and integer
positions $q,q'$, updating the work tape and head position of the represented
configuration with state $s'$, input-head position $ip$, stored value $v'$, and
work-tape head at $q'$ yields the represented configuration with state $s$, input-head
position $ip$, stored value $v'$, and work-tape head at $q$.
\end{lemma}

\begin{definition}[Bounded work-head reachability]
\lean{Complexity.LogProg.Dbl.Rch}
\label{Complexity.LogProg.Dbl.Rch}
\leanok
\uses{Complexity.LogProg.Dbl.St, Complexity.LogProg.Dbl.o, Complexity.LogProg.Dbl.prog, Complexity.LogProg.rrun, Turing.Cfg}
For a Boolean input string $y$, a natural number $W$, and configurations $c,c'$ of the
one-work-tape machine on input $y$, there is a time $T$ such that running the program
from $c$ for $T$ steps reaches $c'$, and at every earlier time $t<T$ the work-head
position lies between $-1$ and $W$, inclusive.
\end{definition}

\begin{lemma}[Transitivity of Bounded Reachability]
\lean{Complexity.LogProg.Dbl.Rch.trans}
\label{Complexity.LogProg.Dbl.Rch.trans}
\leanok
\uses{Complexity.LogProg.Dbl.Rch, Complexity.LogProg.Dbl.St, Complexity.LogProg.rrun_add, Turing.Cfg}
\difficulty{4}
Fix an input word $y$ and a bound $W\in\mathbb{N}$. For the specified machine and
oracle, suppose there is a finite run from configuration $a$ to $b$ and a finite run
from $b$ to $c$, such that at every time before each run's endpoint the work-tape head
lies between $-1$ and $W$, inclusive. Then there is a finite run from $a$ to $c$ with
the same head-position bounds at every time before its endpoint.
\end{lemma}

\begin{lemma}[One-step transition for input movement]
\lean{Complexity.LogProg.Dbl.step_x}
\label{Complexity.LogProg.Dbl.step_x}
\leanok
\uses{Complexity.LogProg.Dbl.K, Complexity.LogProg.Dbl.St, Complexity.LogProg.Dbl.o, Complexity.LogProg.Dbl.prog, Complexity.LogProg.Dbl.tr, Complexity.LogProg.Dbl.xCfg_K, Complexity.LogProg.apply_xAct, Complexity.LogProg.inSym, Complexity.LogProg.rrun, Complexity.LogProg.rrun_one_x, Complexity.LogProg.xAct, Turing.Action.apply, Turing.moveInputPos}
\difficulty{3}
Let $s,s'$ be states, let $ip$ be a valid input-head position on a Boolean input word
$y$, let $v$ be a natural number, and let $q$ be an integer. Suppose that for every
possible work-tape reading, the transition from state $s$ on the symbol read at input
position $ip$ is the action that moves the input head by $mvI$, leaves the work tape
unchanged, and enters state $s'$. Then one step of the doubled-input machine from the
configuration with state $s$, input-head position $ip$, tape storing the binary
expansion of $v$, and work-tape head at $q$ gives the configuration with state $s'$,
input-head position moved by the clamped movement $mvI$, the same tape, and work-tape
head still at $q$.
\end{lemma}

\begin{lemma}[Unary prefix of a pair encoding]
\lean{Complexity.LogProg.Dbl.sym_lt}
\label{Complexity.LogProg.Dbl.sym_lt}
\leanok
\uses{Complexity.LogProg.dbl_replicate, Turing.pairEncode, Turing.pairEncode_eq_dbl}
\difficulty{3}
Let $n$ be a natural number, let $w$ be a Boolean word, and let $y$ be the pair encoding
of the word of $n$ ones and $w$. For every natural number $k<2n$, the symbol of $y$ at
zero-based position $k$ is true.
\end{lemma}

\begin{lemma}[Separator symbol in a unary pair encoding]
\lean{Complexity.LogProg.Dbl.sym_eq}
\label{Complexity.LogProg.Dbl.sym_eq}
\leanok
\uses{Complexity.LogProg.dbl_replicate, Turing.pairEncode, Turing.pairEncode_eq_dbl}
\difficulty{3}
Let $n$ be a natural number, let $w$ be a Boolean word, and let $y$ be the pair encoding
of the word of $n$ ones and $w$. The symbol of $y$ at zero-based position $2n$ is false.
\end{lemma}

\begin{lemma}[Length of an Encoded Unary Pair]
\lean{Complexity.LogProg.Dbl.length_eq}
\label{Complexity.LogProg.Dbl.length_eq}
\leanok
\uses{Complexity.LogProg.dbl_replicate, Turing.pairEncode, Turing.pairEncode_eq_dbl}
\difficulty{3}
Let $n\in\mathbb{N}$ and let $w$ be a Boolean word. If $y$ is the self-delimiting
encoding of the unary word $1^n$ paired with $w$, then $y$ has length $2n+2+|w|$.
\end{lemma}

\begin{lemma}[Bounded Counter Run]
\lean{Complexity.LogProg.Dbl.count}
\label{Complexity.LogProg.Dbl.count}
\leanok
\uses{Complexity.LogProg.Dbl.K, Complexity.LogProg.Dbl.Rch, Complexity.LogProg.Dbl.Rch.trans, Complexity.LogProg.Dbl.length_eq, Complexity.LogProg.Dbl.o, Complexity.LogProg.Dbl.prog, Complexity.LogProg.Dbl.regCfg_K, Complexity.LogProg.Dbl.step_x, Complexity.LogProg.Dbl.sym_lt, Complexity.LogProg.Dbl.tr, Complexity.LogProg.inSym_succ, Complexity.LogProg.inc_run, Complexity.LogProg.length_bits_mono, Complexity.LogProg.moveInputPos_pos_val, Complexity.LogProg.regCfg, Complexity.LogProg.rrun, Complexity.LogProg.rrun_add, Complexity.LogProg.rrun_zero, Turing.pairEncode}
\difficulty{6}
Let $n\in\mathbb{N}$ and let $w$ be a Boolean word. If the input is the self-delimiting
encoding of the unary word $1^n$ paired with $w$, then for every $j\le 2n$ the specified
machine has a finite run from its counter state with input head at position $1$, counter
value $0$, and work-tape head at $0$, to its counter state with input head at position
$1+j$, counter value $j$, and work-tape head at $0$. Before the run ends, the work-tape
head stays between $-1$ and $\lvert\operatorname{bits}(2n)\rvert$, inclusive.
\end{lemma}

\begin{lemma}[Initial Configuration Representation]
\lean{Complexity.LogProg.Dbl.init_eq}
\label{Complexity.LogProg.Dbl.init_eq}
\leanok
\uses{Complexity.LogProg.Dbl.K, Turing.Cfg.init}
\difficulty{3}
For every Boolean input word $y$, the initial configuration of the specified one-tape
machine in state $vU0$ on input $y$ equals the configuration with state $vU0$,
input-head position $1$, the binary expansion of $0$ stored on its work tape beginning
at cell zero, work-tape head at position $0$, and empty output.
\end{lemma}

\begin{lemma}[Validity of the doubled-word run]
\lean{Complexity.LogProg.Dbl.run_valid}
\label{Complexity.LogProg.Dbl.run_valid}
\leanok
\uses{Complexity.LogProg.Canon, Complexity.LogProg.Dbl.K, Complexity.LogProg.Dbl.Rch, Complexity.LogProg.Dbl.Rch.trans, Complexity.LogProg.Dbl.count, Complexity.LogProg.Dbl.length_eq, Complexity.LogProg.Dbl.o, Complexity.LogProg.Dbl.prog, Complexity.LogProg.Dbl.step_x, Complexity.LogProg.Dbl.sym_eq, Complexity.LogProg.Dbl.tr, Complexity.LogProg.Dbl.xCfg_K, Complexity.LogProg.inSym_succ, Complexity.LogProg.jeqPlain_run, Complexity.LogProg.rewind_x, Complexity.LogProg.rrun_zero, Complexity.LogProg.valPlain_run, Turing.moveInputPos_zero, Turing.pairEncode}
\difficulty{5}
Let $n$ be a natural number and let $w$ be a Boolean word whose last bit is true
whenever $w$ is nonempty. Suppose the input is the pair encoding of the word of $n$ ones
and $w$. Then the doubled-input machine has a run from its initial configuration, with
input head at position $1$, empty output, and work-tape head at $0$, to a configuration
with input head at position $1$, work-tape head at $0$, and state yes exactly when $w$
is the binary expansion of $2n$, and no otherwise. Throughout the run, at every time
before its endpoint, the work-tape head position lies between $-1$ and the length of the
binary expansion of $2n$, inclusive.
\end{lemma}

\begin{lemma}[Halting return step]
\lean{Complexity.LogProg.Dbl.ret_step}
\label{Complexity.LogProg.Dbl.ret_step}
\leanok
\uses{Complexity.LogProg.Dbl.K, Complexity.LogProg.Dbl.o, Complexity.LogProg.Dbl.prog, Complexity.LogProg.Dbl.tr, Complexity.LogProg.retAct, Complexity.LogProg.rrun, Complexity.LogProg.rrun_one, Complexity.LogProg.rstep_noncall, Turing.Action.apply, Turing.MultiTapeTM.step}
\difficulty{4}
For every Boolean $b$, input-head position $ip$ on a fixed input word $y$, and natural
number $v$, one step of the doubled-input machine from the configuration in state yes
when $b$ is true and no when $b$ is false, with work tape storing the binary word for
$v$ and its head at $0$, reaches a halting configuration. Its output is the one-symbol
word $[b]$, and its work-tape head remains at $0$.
\end{lemma}

\begin{definition}[Machine with no work tapes]
\lean{Complexity.LogProg.Dbl.nilTM}
\label{Complexity.LogProg.Dbl.nilTM}
\leanok
\uses{Turing.MultiTapeTM}
The machine has no work tapes, a single state, and Boolean input and output alphabets.
It begins in that state; on every transition it leaves the input head and output
unchanged, writes nothing, and halts.
\end{definition}

\begin{theorem}[Logarithmic-space membership]
\lean{Complexity.LogProg.dblLang_mem}
\label{Complexity.LogProg.dblLang_mem}
\leanok
\uses{Complexity.LOGSPACE, Complexity.LogProg.Dbl.K, Complexity.LogProg.Dbl.init_eq, Complexity.LogProg.Dbl.length_eq, Complexity.LogProg.Dbl.nilTM, Complexity.LogProg.Dbl.o, Complexity.LogProg.Dbl.prog, Complexity.LogProg.Dbl.ret_step, Complexity.LogProg.Dbl.run_valid, Complexity.LogProg.Dbl.xCfg_K, Complexity.LogProg.ValidPlain, Complexity.LogProg.canon_bits, Complexity.LogProg.compileFinTM, Complexity.LogProg.compile_space, Complexity.LogProg.dblLang, Complexity.LogProg.indicator_eq_decide, Complexity.LogProg.length_bits_le_log, Complexity.LogProg.length_bits_mono, Complexity.LogProg.rrun, Complexity.LogProg.rrun_add, Complexity.LogProg.valPlain_run, Complexity.logSpace, Turing.Cfg.init, Turing.MultiTapeTM.indicator, Turing.pairEncode_injective}
\difficulty{7}
The language of Boolean words of the form $\langle 1^n,\operatorname{bits}(2n)\rangle$,
with $n\in\mathbb{N}$ and the pairing formed by doubling each bit of the first word,
appending the separator $01$, and then the second word, belongs to logarithmic space.
Here logarithmic space means that some finite-state multi-tape Turing machine decides
membership while using at most a constant multiple of $\lfloor\log_2 N\rfloor+1$
work-tape cells on inputs of length $N$.
\end{theorem}
